% GENERATED FILE. Edit canonical lessons, then run npm run generate:books.
% canonical-source-hash: fnv1a64:e7d4ba3cbb7a3eda
% canonical-lessons: GE-C353-betrieb, GE-C353-unternehmen, GE-C353-abteilung, GE-C353-vorgesetzte, GE-C353-leiten, GE-R353-review-346-353

\chapter{Im Betrieb}
\label{ch:ge-im-betrieb}
\hlchaptermodality{353}

\begin{chapteropening}
\textbf{By the end of this chapter:} \emph{I can talk about a firm: the company, its departments, my superior and who runs what.}

Complete the last lesson of chapter 353: I can talk about a firm: the company, its departments, my superior and who runs what.
\end{chapteropening}

\section[der Betrieb]{der Betrieb — the firm, the business; operation}

\label{lesson:GE-C353-betrieb}



\begin{quote}
Before the new one: say the German for the trial period, then the German for to train in.
\end{quote}

\subsection*{You'll want to know: der Betrieb}
\textbf{der Betrieb} — \textquotedblleft{}the firm, the business; operation\textquotedblright{}.

\begin{quote}
Mein Onkel hat einen kleinen Betrieb mit acht Mitarbeitern.
\end{quote}

\begin{quote}
\emph{My uncle has a small business with eight employees.}
\end{quote}

Plural: \emph{die Betriebe}. \emph{In Betrieb} — in operation; \emph{außer Betrieb} — out of order.

\begin{cousinweb}
\emph{Betrieb} is the noun of \emph{betreiben} \textquotedblleft{}to run, to operate\textquotedblright{}, built with \emph{be-} on \emph{treiben} \textquotedblleft{}to drive\textquotedblright{}, inherited from Proto-Germanic \emph{drībaną} (EtymDB), a cousin of English \emph{drive}.
\end{cousinweb}

\begin{grammarlens}[title={What you've built}]
One more word for telling, explaining and describing at B1.
\end{grammarlens}

\subsection*{Your turn}
\begin{itemize}
\raggedright
  \item \emph{Say it:} \emph{der Betrieb}
  \item \emph{Say it:} \emph{der Betrieb}, once more
  \item \emph{Say it:} a sentence of your own with \emph{der Betrieb}
  \item \emph{Recall:} say the German for the trial period, then the German for to train in, then say \emph{der Betrieb} again
\end{itemize}

\begin{tcolorbox}[breakable,colback=teal!4,colframe=teal!35!black,title={Before you move on}]
What does der Betrieb mean? (\textquotedblleft{}The firm, the business; operation\textquotedblright{}.) Say it once more.
\end{tcolorbox}
\section[das Unternehmen]{das Unternehmen — the company, the enterprise}

\label{lesson:GE-C353-unternehmen}



\begin{quote}
Before the new one: say the German for to train in, then the German for the firm.
\end{quote}

\subsection*{You'll want to know: das Unternehmen}
\textbf{das Unternehmen} — \textquotedblleft{}the company, the enterprise\textquotedblright{}.

\begin{quote}
Das Unternehmen verkauft Möbel in vielen Ländern.
\end{quote}

\begin{quote}
\emph{The company sells furniture in many countries.}
\end{quote}

Plural: \emph{die Unternehmen}. Also the verb \emph{unternehmen} — to undertake, to do: \emph{Was unternehmen wir heute?}

\begin{cousinweb}
\emph{Unternehmen} is the verb \emph{unternehmen} \textquotedblleft{}to undertake\textquotedblright{} made into a noun; the verb is built from \emph{unter} and \emph{nehmen} \textquotedblleft{}to take\textquotedblright{}, Old High German \emph{neman}, Proto-Germanic \emph{nemaną} (EtymDB).
\end{cousinweb}

\begin{grammarlens}[title={What you've built}]
One more word for telling, explaining and describing at B1.
\end{grammarlens}

\subsection*{Your turn}
\begin{itemize}
\raggedright
  \item \emph{Say it:} \emph{das Unternehmen}
  \item \emph{Say it:} \emph{das Unternehmen}, once more
  \item \emph{Say it:} a sentence of your own with \emph{das Unternehmen}
  \item \emph{Recall:} say the German for to train in, then the German for the firm, then say \emph{das Unternehmen} again
\end{itemize}

\begin{tcolorbox}[breakable,colback=teal!4,colframe=teal!35!black,title={Before you move on}]
What does das Unternehmen mean? (\textquotedblleft{}The company, the enterprise\textquotedblright{}.) Say it once more.
\end{tcolorbox}
\section[die Abteilung]{die Abteilung — the department}

\label{lesson:GE-C353-abteilung}



\begin{quote}
Before the new one: say the German for the firm, then the German for the company.
\end{quote}

\subsection*{You'll want to know: die Abteilung}
\textbf{die Abteilung} — \textquotedblleft{}the department\textquotedblright{}.

\begin{quote}
Sie arbeitet in der Abteilung für Marketing.
\end{quote}

\begin{quote}
\emph{She works in the marketing department.}
\end{quote}

Plural: \emph{die Abteilungen}. In a hospital, a ward: \emph{die Kinderabteilung}.

\begin{cousinweb}
\emph{Abteilung} is the noun of \emph{abteilen} \textquotedblleft{}to divide off\textquotedblright{}, built on \emph{Teil} \textquotedblleft{}part\textquotedblright{}, Old High German \emph{teil}, Proto-Germanic \emph{dailą} (EtymDB): one part of the whole.
\end{cousinweb}

\begin{grammarlens}[title={What you've built}]
One more word for telling, explaining and describing at B1.
\end{grammarlens}

\subsection*{Your turn}
\begin{itemize}
\raggedright
  \item \emph{Say it:} \emph{die Abteilung}
  \item \emph{Say it:} \emph{die Abteilung}, once more
  \item \emph{Say it:} a sentence of your own with \emph{die Abteilung}
  \item \emph{Recall:} say the German for the firm, then the German for the company, then say \emph{die Abteilung} again
\end{itemize}

\begin{tcolorbox}[breakable,colback=teal!4,colframe=teal!35!black,title={Before you move on}]
What does die Abteilung mean? (\textquotedblleft{}The department\textquotedblright{}.) Say it once more.
\end{tcolorbox}
\section[der Vorgesetzte]{der Vorgesetzte — the superior, the boss}

\label{lesson:GE-C353-vorgesetzte}



\begin{quote}
Before the new one: say the German for the company, then the German for the department.
\end{quote}

\subsection*{You'll want to know: der Vorgesetzte}
\textbf{der Vorgesetzte} — \textquotedblleft{}the superior, the boss\textquotedblright{}.

\begin{quote}
Mein Vorgesetzter ist sehr freundlich.
\end{quote}

\begin{quote}
\emph{My superior is very friendly.}
\end{quote}

Declined like an adjective: \emph{der Vorgesetzte}, \emph{ein Vorgesetzter}, and \emph{die Vorgesetzte} for a woman.

\begin{cousinweb}
\emph{Vorgesetzte} is the participle of \emph{vorsetzen} \textquotedblleft{}to set before, to set over\textquotedblright{}, built on \emph{setzen} \textquotedblleft{}to set\textquotedblright{}, Proto-Germanic \emph{satjaną} (EtymDB), made into a noun: the person set over you.
\end{cousinweb}

\begin{grammarlens}[title={What you've built}]
One more word for telling, explaining and describing at B1.
\end{grammarlens}

\subsection*{Your turn}
\begin{itemize}
\raggedright
  \item \emph{Say it:} \emph{der Vorgesetzte}
  \item \emph{Say it:} \emph{der Vorgesetzte}, once more
  \item \emph{Say it:} a sentence of your own with \emph{der Vorgesetzte}
  \item \emph{Recall:} say the German for the company, then the German for the department, then say \emph{der Vorgesetzte} again
\end{itemize}

\begin{tcolorbox}[breakable,colback=teal!4,colframe=teal!35!black,title={Before you move on}]
What does der Vorgesetzte mean? (\textquotedblleft{}The superior, the boss\textquotedblright{}.) Say it once more.
\end{tcolorbox}
\section[leiten]{leiten — to run, to lead (a team, a firm)}

\label{lesson:GE-C353-leiten}



\begin{quote}
Before the new one: say the German for the department, then the German for the superior.
\end{quote}

\subsection*{You'll want to know: leiten}
\textbf{leiten} — \textquotedblleft{}to run, to lead (a team, a firm)\textquotedblright{}.

\begin{quote}
Meine Mutter leitet eine Schule mit vierhundert Kindern.
\end{quote}

\begin{quote}
\emph{My mother runs a school with four hundred children.}
\end{quote}

\emph{Die Leitung} — the management. \emph{Der Leiter, die Leiterin} — the head.

\begin{cousinweb}
\emph{leiten} is inherited from Middle High German \emph{leiten}, Old High German \emph{leiten} and Proto-Germanic \emph{laidijaną} \textquotedblleft{}to lead\textquotedblright{}, built on \emph{līþaną} \textquotedblleft{}to go\textquotedblright{} (EtymDB); English \emph{lead} is a cousin.
\end{cousinweb}

\begin{grammarlens}[title={What you've built}]
One more word for telling, explaining and describing at B1.
\end{grammarlens}

\subsection*{Your turn}
\begin{itemize}
\raggedright
  \item \emph{Say it:} \emph{leiten}
  \item \emph{Say it:} \emph{leiten}, once more
  \item \emph{Say it:} a sentence of your own with \emph{leiten}
  \item \emph{Recall:} say the German for the department, then the German for the superior, then say \emph{leiten} again
\end{itemize}

\begin{tcolorbox}[breakable,colback=teal!4,colframe=teal!35!black,title={Before you move on}]
What does leiten mean? (\textquotedblleft{}To run, to lead (a team, a firm)\textquotedblright{}.) Say it once more.
\end{tcolorbox}
\section[(dialogue)]{Words from chapters 346–353 — a review}

\label{lesson:GE-R353-review-346-353}



\begin{quote}
Say the German for the purpose and the intention.
\end{quote}

\subsection*{Your turn}
Say these aloud:

\begin{itemize}
\raggedright
  \item the purpose, the intention, in order that, on purpose, the occasion
  \item the chance, by chance, to prevent, to avoid, why
  \item the responsibility, responsible, to justify, the excuse, that is why
  \item the job advertisement, the job interview, the ability, qualified, to master
  \item the employer, the employee, to hire, to employ, fixed-term
  \item part-time, full-time, the shift, the hour of overtime, the end of the working day
  \item to lay off, unemployed, unemployment, the trial period, to train in
  \item the firm, the company, the department, the superior, to run
\end{itemize}

\begin{tcolorbox}[breakable,colback=teal!4,colframe=teal!35!black,title={Before you move on}]
The purpose? (\textbf{der Zweck}.) The intention? (\textbf{die Absicht}.) In order that? (\textbf{damit}.) On purpose? (\textbf{absichtlich}.) The occasion? (\textbf{der Anlass}.) The chance? (\textbf{der Zufall}.) By chance? (\textbf{zufällig}.) To prevent? (\textbf{verhindern}.) To avoid? (\textbf{vermeiden}.) Why? (\textbf{weshalb}.) The responsibility? (\textbf{die Verantwortung}.) Responsible? (\textbf{verantwortlich}.) To justify? (\textbf{rechtfertigen}.) The excuse? (\textbf{die Ausrede}.) That is why? (\textbf{deswegen}.) The job advertisement? (\textbf{die Stellenanzeige}.) The job interview? (\textbf{das Vorstellungsgespräch}.) The ability? (\textbf{die Fähigkeit}.) Qualified? (\textbf{qualifiziert}.) To master? (\textbf{beherrschen}.) The employer? (\textbf{der Arbeitgeber}.) The employee? (\textbf{der Arbeitnehmer}.) To hire? (\textbf{einstellen}.) To employ? (\textbf{beschäftigen}.) Fixed-term? (\textbf{befristet}.) Part-time? (\textbf{die Teilzeit}.) Full-time? (\textbf{die Vollzeit}.) The shift? (\textbf{die Schicht}.) The hour of overtime? (\textbf{die Überstunde}.) The end of the working day? (\textbf{der Feierabend}.) To lay off? (\textbf{entlassen}.) Unemployed? (\textbf{arbeitslos}.) Unemployment? (\textbf{die Arbeitslosigkeit}.) The trial period? (\textbf{die Probezeit}.) To train in? (\textbf{einarbeiten}.) The firm? (\textbf{der Betrieb}.) The company? (\textbf{das Unternehmen}.) The department? (\textbf{die Abteilung}.) The superior? (\textbf{der Vorgesetzte}.) To run? (\textbf{leiten}.)
\end{tcolorbox}
